It remains to explain that the top Segre class of $L^{[k]}$ is positive. We use \eqref{lehnc} and we change variables
\begin{gather*}
z=t(1+2t)^2\qquad \text{so that}\quad \mathrm dz=(1+2t)(1+6t)\,{\mathrm dt}.
\end{gather*}
Then
\begin{align}
\int_{X^{[k]}} s\big(L^{[k]}\big)&=\operatorname{Res}_{z=0} A_1^{2h-\ell^2} A_2^2 A_3^{\ell} A_4^{-1} \frac{\mathrm dz}{z^{k+1}}\nonumber
\\
\nonumber&=
\operatorname{Res}_{t=0} 2^{-\ell-2} (1+2t)^{h-\frac{\ell^2}{2}-2k-\ell+\frac{3}{2}} (1+6t)^{-\frac{1}{2}} \left(\sqrt{1+2t}+\sqrt{1+6t}\right)^{\ell+2}\frac{\mathrm dt}{t^{k+1}}
\\
&= 2^{-\ell-2} \operatorname{Coeff}_{t^k} (1\!+2t)^{h-\frac{\ell^2}{2}-2k-\ell+\frac{3}{2}} (1\!+6t)^{-\frac{1}{2}} \left(\sqrt{1+2t}\!+\sqrt{1+6t}\right)^{\ell+2}\!.
\label{cvvv}
\end{align}
